\chapter{Related Work}
\label{chap:relatedwork}
\pagestyle{plain}

\section{Learned Cardinality Estimation}

Machine learning-based approaches for replacing traditional database components have gained significant attention, particularly in the context of cardinality estimation. A foundational contribution in this area is the Multi-Set Convolutional Network (MSCN) proposed by Kipf et al.~\cite{DBLP:conf/cidr/KipfKRLBK19}, which reformulates cardinality estimation as a supervised learning problem. Unlike conventional estimators that rely on simplifying assumptions such as attribute independence, MSCN represents queries as sets of tables, joins, and predicates, enabling the model to capture complex correlations across query components.

MSCN demonstrated that learned models could substantially outperform traditional histogram-based estimators on complex analytical query workloads.
Its permutation-invariant architecture and use of sampled bitmap representations enabled improved modeling of complex correlations across joins and predicates,
establishing MSCN as a foundational baseline for subsequent learned cardinality estimation research.

Despite these advantages, MSCN fundamentally depends on execution-based labeling. Each training instance requires executing a query on a database to obtain its true cardinality. This introduces substantial computational overhead and limits scalability, particularly when large and diverse workloads are required. Moreover, adapting to new schemas or workloads necessitates additional query execution, further increasing the cost. These limitations highlight a central challenge in learned cardinality estimation: the dependency on labeled data generated through expensive query execution.

\section{Workload Characteristics and Model Generalization}

Beyond model architecture, the effectiveness of learned cardinality estimators is strongly influenced by the characteristics of the training workload. Wang et al.~\cite{DBLP:journals/pvldb/WangQWWZ21} demonstrate that learned models are highly sensitive to factors such as predicate types, selectivity, join complexity, and query structure.

Their findings show that models trained on narrow or biased workloads often struggle to generalize to unseen query patterns. In particular, learned estimators may experience substantial performance degradation when evaluated on workloads containing predicate types, selectivity ranges, or join structures that differ from those observed during training.

The issue of workload shift further complicates deployment. In real-world systems, query distributions evolve over time, leading to discrepancies between training and inference workloads. Under such conditions, learned models can exhibit significant decreases in estimation accuracy. These observations emphasize the importance of constructing training datasets that are both diverse and representative of real-world query patterns.

Recent work on robust cardinality estimation under changing and evolving workloads~\cite{DBLP:journals/pvldb/NegiWKTMMKA23} further highlights the importance of adaptability, reinforcing the need for scalable data generation approaches capable of capturing a broad range of query characteristics.

\section{LLM-based SQL Workload Generation}

The emergence of large language models (LLMs) has introduced new possibilities for SQL workload generation. The SQLStorm framework~\cite{DBLP:journals/pvldb/SchmidtLBN25} demonstrates that LLMs can generate large-scale, diverse, and realistic SQL queries using schema-aware prompting techniques.

Unlike traditional template-based approaches, which are inherently limited in structural diversity, LLM-based methods can generate queries with varying levels of complexity, including different join structures, predicate combinations, and nested queries. This flexibility enables the creation of more realistic workloads that better reflect real-world usage patterns and expose learned models to broader query distributions.

Recent advances in text-to-SQL and structured reasoning further support the broader applicability of LLMs for SQL generation tasks. Systems such as CHASE-SQL~\cite{DBLP:conf/iclr/PourrezaL0CTKGS25} and TAPEX~\cite{DBLP:conf/iclr/LiuCGZLCL22} demonstrate that neural models are increasingly capable of understanding relational schemas, table structures, and complex query semantics. Evaluation studies by Rahaman et al.~\cite{DBLP:conf/edbt/RahmanZMCP25} additionally highlight both the strengths and limitations of contemporary LLMs in structured query generation settings.

However, existing LLM-based approaches primarily focus on query generation rather than labeled data creation. Generated queries must still be executed against the database in order to obtain ground-truth cardinalities, thereby retaining the dependency on execution-based labeling and limiting scalability for learned cardinality estimation pipelines.

Furthermore, existing LLM-based approaches evaluate their models on externally curated benchmark workloads rather than on LLM-generated queries. In contrast, the approach presented in this thesis supports a self-contained pipeline in which LLM-generated queries are used for training and validation in LLM-based runs, while controlled comparisons may additionally use synthetic workloads for reproducibility.

\section{Generative AI for Training Database Components}

Recent work on integrating generative AI into database systems has sought to address the limitations of manual and execution-based data generation. Davitkova and Michel~\cite{DBLP:journals/corr/abs-2512-20271} propose a framework that leverages generative models to automate the training of learned database components.

Their framework introduces a structured pipeline that integrates query generation, validation, execution, and model training. By automating query generation, it reduces reliance on manually collected query logs and enables the creation of diverse training datasets. Techniques such as schema-aware prompting and iterative refinement are further used to improve the quality and realism of generated queries.

Despite these advancements, generated queries must still be executed to obtain ground-truth cardinalities for supervised training. Consequently, query execution continues to represent a significant computational cost within learned cardinality estimation pipelines.

\section{Active Learning and Uncertainty Quantification}

Active learning provides a complementary approach to improving data efficiency in supervised learning systems. A recent survey on deep active learning~\cite{DBLP:journals/tnn/LiWCJDO25} describes how active learning methods aim to reduce labeling requirements by selecting informative samples for annotation. A core mechanism for this selection is uncertainty quantification, which allows a model to identify the inputs it is least confident about.

A theoretically grounded method for estimating neural network uncertainty is Monte Carlo (MC) dropout. The regularizing and averaging properties of dropout were fundamentally established by Baldi and Sadowski~\cite{DBLP:journals/ai/BaldiS14}, who demonstrated that dropout acts as a computationally efficient approximation of an exponentially large ensemble of subnetworks. By retaining dropout during inference and performing multiple stochastic forward passes, models can approximate predictive uncertainty, with variation across predictions often used as an estimate of epistemic uncertainty~\cite{DBLP:conf/icml/GalG16}\cite{hossain2026uatliteinferencetimeuncertaintyawareattention}.

While MC dropout provides useful approximate uncertainty estimates, its repeated forward passes introduce substantial computational overhead\cite{hossain2026uatliteinferencetimeuncertaintyawareattention}. This issue is particularly problematic when querying large unlabeled pools during active learning. To address this, Srinivasan~\cite{Srinivasan2025Efficient} recently explored model-splitting techniques that partition the network into a deterministic body and a stochastic head. This architectural optimization achieves up to a 25\--33x speedup in inference, transforming MC dropout from a theoretically appealing concept into a highly practical tool for rapid uncertainty quantification.

In the context of learned cardinality estimation, active learning can be used to minimize the number of queries that need to be executed for expensive labeling. Prior work on active learning for ML-enhanced database systems has shown that additional query labels can be collected cost-effectively to improve deployed model performance~\cite{DBLP:conf/sigmod/MaDDS20}. By prioritizing queries where the model is uncertain or performs poorly, it is possible to achieve high estimation accuracy with fewer labeled samples.

In addition to uncertainty-based strategies, diversity-based sampling plays an important role in ensuring that selected queries cover a wide range of patterns. This is particularly relevant when working with synthetically generated data, where redundancy may otherwise reduce training effectiveness. Combining uncertainty and diversity criteria allows for the construction of efficient and representative training datasets.

While active learning is well established within machine learning research, its use in learned database systems has received comparatively limited research attention. Integrating active learning with LLM-based query generation provides a practical and scalable mechanism for improving labeling efficiency in learned cardinality estimation systems.

\section{Data Diversity and Distribution Alignment}

Ensuring that synthetic workloads reflect real-world query distributions is important for training robust learned database systems. The PAARS framework~\cite{DBLP:journals/corr/abs-2503-24228} explores persona-based generation strategies to encourage greater variation in generated data distributions and improve workload diversity.

A key contribution of this work is the use of KL divergence to measure similarity between generated and real-world data distributions. This provides a quantitative mechanism for evaluating whether synthetic workloads capture important structural characteristics of observed query patterns. In the context of SQL workloads, such metrics can be applied to distributions over query features including join counts, predicate types, and selectivity ranges.

Distribution alignment is particularly important for learned cardinality estimation because models trained on unrealistic synthetic workloads may fail to generalize effectively to real query patterns. Incorporating distributional similarity metrics into workload generation therefore helps improve both dataset diversity and representativeness.

\section{Positioning of the Proposed Approach}

The existing literature demonstrates significant recent progress in learned database systems and LLM-based data generation. Models such as MSCN~\cite{DBLP:conf/cidr/KipfKRLBK19} highlight the potential of learned approaches for cardinality estimation, while studies on workload characteristics~\cite{DBLP:journals/pvldb/WangQWWZ21} emphasize the importance of diverse and representative training data. LLM-based frameworks such as SQLStorm~\cite{DBLP:journals/pvldb/SchmidtLBN25} and generative pipelines~\cite{DBLP:journals/corr/abs-2512-20271} further demonstrate the practical feasibility of scalable SQL workload generation.

However, a common limitation across these approaches is the continued reliance on execution-based labeling. Although query generation has become increasingly automated, obtaining ground-truth labels still requires executing queries on a database, introducing substantial computational overhead.

The approach proposed in this thesis addresses this limitation by leveraging LLMs to generate diverse SQL workloads while using active learning to minimize the number of expensive database executions required for labeling. Rather than eliminating execution-based supervision entirely, the framework is designed to reduce labeling cost by selecting informative queries for execution and training. Furthermore, by incorporating active learning~\cite{DBLP:journals/tnn/LiWCJDO25} and diversity-aware generation strategies~\cite{DBLP:journals/corr/abs-2503-24228}, the proposed method aims to improve both data efficiency and workload representativeness.

Within the proposed framework, the LLM acts as a workload generator, while the trained supervised model serves as the final learned database component during inference. This separation enables scalable training workflows and supports more efficient development of learned cardinality estimation systems.

The primary contribution of this work does not lie in the introduction of a novel neural architecture for cardinality estimation, but rather in the design and evaluation of an automated training framework that integrates LLM-based workload generation with uncertainty-guided active learning. The proposed approach specifically addresses the practical challenges associated with execution-based labeling by reducing the number of required database executions while maintaining workload diversity and estimation quality. In this context, the thesis investigates how automated workload generation and active learning strategies can improve the scalability and data efficiency of training learned cardinality estimation systems.